% SEGMENT OLP-0035-S01
% Part:sets-functions-relations
% Chapter: size-of-sets
% Section: equinumerous-sets

\documentclass[../../../include/open-logic-section]{subfiles}

\begin{document}

\olfileid{sfr}{siz}{equ}

\olsection{எண்ணொத்த தன்மை}

கணங்களின் “அளவு” பற்றிய ஓர் உள்ளுணர்வுக் கருத்து நம்மிடம் உள்ளது;
அது முடிவுறு கணங்களுக்கு நன்றாகச் செயல்படுகிறது.
ஆனால் முடிவுறாத கணங்களின் நிலை என்ன?
\emph{எந்த} அளவைக் கொண்ட இரண்டு கணங்களின் அளவுகளையும் ஒப்பிடுவதற்கான
முறையான ஒரு வழியை உருவாக்க வேண்டுமானால்,
கணங்கள் எப்போது ஒரே அளவைக் கொண்டுள்ளன என்பதை வரையறுப்பதிலிருந்து
தொடங்குவது நல்லது. ஃப்ரேகே கூறுவது இதோ:
\begin{quote}
ஒரு பணியாளர், மேசையில் தட்டுகளின் எண்ணிக்கைக்குச் சரியாகச் சமமான
எண்ணிக்கையில் கத்திகளை வைத்திருக்கிறாரா என்பதை உறுதிசெய்ய விரும்பினால்,
அவற்றில் எதையும் அவர் எண்ண வேண்டியதில்லை.
ஒவ்வொரு தட்டின் வலப்புறத்திலும் ஒரு கத்தியை வைத்து,
மேசையின் ஒவ்வொரு கத்தியும் ஏதாவது ஒரு தட்டின் வலப்புறத்தில்
இருக்குமாறு செய்தால் போதும்.
இவ்வாறு தட்டுகளும் கத்திகளும் ஒன்றுடன் ஒன்று தனித்தனியாகத்
தொடர்புபடுத்தப்படுகின்றன; அதுவும் அதே இடம்சார் தொடர்பின் வழியாகவே.
\citep[\S70]{Frege1884}
\end{quote}
இந்தப் பத்தியின் உட்பார்வையை ஒரு முறையான வரையறையின் மூலம்
வெளிப்படுத்தலாம்:

% SEGMENT OLP-0035-S02
\begin{defn}\ollabel{comparisondef}
!!a{bijection} $f \colon A \to B$ இருக்கும்போது,
அப்போதும் அப்போது மட்டுமே, $A$ ஆனது $B$ உடன்
\emph{எண்ணொத்தது} எனப்படும்; இதனை $\cardeq{A}{B}$ என எழுதுகிறோம்.
\end{defn}

% SEGMENT OLP-0035-S03
\begin{prop}\ollabel{equinumerosityisequi}
எண்ணொத்த தன்மை ஒரு சமானத் தொடர்பாகும்.
\end{prop}

% SEGMENT OLP-0035-S04
\begin{proof}
எண்ணொத்த தன்மை தற்சுட்டு, சமச்சீர், கடப்பு ஆகிய பண்புகளைக்
கொண்டுள்ளது எனக் காட்ட வேண்டும்.
$A, B$, $C$ ஆகியவை கணங்களாக இருக்கட்டும்.

\emph{தற்சுட்டுப் பண்பு.}
அனைத்து $x \in A$ க்கும் $\Id{A} (x) = x$ என அமையும்
ஒன்றாதல் சார்பு $\Id{A} \colon A \to A$ என்பது !!a{bijection}.
எனவே $\cardeq{A}{A}$.

\emph{சமச்சீர்ப் பண்பு.}
$\cardeq{A}{B}$ எனக் கொள்வோம்;
அதாவது !!a{bijection} $f\colon A \to B$ உள்ளது.
$f$ என்பது !!{bijective} பண்புடையதால்,
அதன் நேர்மாறு $f^{-1}$ உள்ளது; அதுவும் !!{bijective} பண்புடையது.
எனவே $f^{-1}\colon B \to A$ என்பது !!a{bijection};
ஆகவே $\cardeq{B}{A}$.

\emph{கடப்புப் பண்பு.}
$\cardeq{A}{B}$, $\cardeq{B}{C}$ எனக் கொள்வோம்.
அதாவது, !!{bijection}s $f\colon A \to B$,
$g\colon B \to
C$ ஆகியவை உள்ளன.
அப்போது சேர்ப்பு $\comp{f}{g}\colon A \to C$ என்பது
!!{bijective} பண்புடையது; ஆகவே $\cardeq{A}{C}$.
\end{proof}

% SEGMENT OLP-0035-S05
\begin{prop}
$\cardeq{A}{B}$ எனில், $A$ என்பது !!{enumerable} ஆக இருக்கும்போது,
அப்போதும் அப்போது மட்டுமே, $B$ உம் அவ்வாறு இருக்கும்.
\end{prop}

% SEGMENT OLP-0035-S06
\begin{editorial}
பின்வரும் நிறுவல் \olref[enm]{defn:enumerable} ஐப் பயன்படுத்துகிறது,
\olref[enm]{sec} சேர்க்கப்பட்டிருந்தால்;
இல்லையெனில் \olref[enm-alt]{defn:enumerable} ஐப் பயன்படுத்துகிறது.
\end{editorial}

% SEGMENT OLP-0035-S07
\begin{proof}
$\cardeq{A}{B}$ எனக் கொள்வோம்.
எனவே ஏதோ ஒரு !!{bijection} $f \colon A
\to B$ உள்ளது; மேலும் $A$ என்பது !!{enumerable} எனக் கொள்வோம்.
% SOURCE CORRECTION OLSIZ-006; both branches use f; see translation/size-notes.tex.
\oliflabeldef{sfr:siz:enm:defn:enumerable}{
  அப்போது $A = \emptyset$ அல்லது
  !!a{surjective} சார்பு $g\colon \PosInt \to A$ உள்ளது.
  $A = \emptyset$ எனில், $B = \emptyset$ உம் ஆகும்.
  ஏனெனில் அவ்வாறு இல்லையெனில், ஏதாவது ஓர் !!a{element}
  $y \in B$ இருக்கும்; ஆனால் $f(x) = y$ என அமையும்
  $x \in A$ இருக்காது.
  மறுபுறம், $g\colon \PosInt \to A$ என்பது !!{surjective}
  பண்புடையது எனில்,
  $\comp{g}{f} \colon \PosInt \to B$ என்பதும் !!{surjective}
  பண்புடையது.
  இதைக் காண, $y \in B$ எனக் கொள்க.
  $f$ என்பது !!{surjective} பண்புடையதால்,
  $f(x) = y$ என அமையும் ஒரு $x \in A$ உள்ளது.
  $g$ என்பது !!{surjective} பண்புடையதால்,
  $g(n) =
  x$ என அமையும் ஒரு $n \in \PosInt$ உள்ளது. எனவே,
\[
(\comp{g}{f})(n) = f(g(n)) = f(x) = y
\]
  ஆகவே $\comp{g}{f}$ என்பது !!{surjective} பண்புடையது.
  $\comp{g}{f}$ என்பது $B$ இன் ஒரு பட்டியலாக்கம்;
  எனவே $B$ என்பது !!{enumerable}.}
{
  அப்போது $A = \emptyset$ அல்லது,
  !!a{bijection} $g$ ஒன்று உள்ளது.
  அதன் வீச்சகம் $A$; அதன் சார்பகம் $\Nat$ அல்லது இயல் எண்களின்
  ஒரு தொடக்கத் தொடராகும்.
  $A = \emptyset$ எனில், $B =
  \emptyset$ உம் ஆகும்.
  ஏனெனில் அவ்வாறு இல்லையெனில், $f(x) = y$ என அமையும்
  $x \in A$ இல்லாத ஏதாவது ஒரு $y \in B$ இருக்கும்.
  ஆகவே நமது !!{bijection} $g$ இருப்பதாகக் கொள்வோம்.
  அப்போது $\comp{g}{f}$ என்பது !!a{bijection}.
  அதன் வீச்சகம் $B$; அதன் சார்பகம் $g$ இன் சார்பகமே
  (அதாவது $\Nat$ அல்லது அதன் ஒரு தொடக்கத் துண்டு).
  எனவே $B$ என்பது !!{enumerable}.}

$B$ என்பது !!{enumerable} எனில்,
$f$ என்பதற்குப் பதிலாக !!{bijection}
$f^{-1}\colon B \to A$ ஐப் பயன்படுத்தி அதே வாதத்தை மீண்டும் செய்வதால்,
$A$ என்பது !!{enumerable} எனப் பெறுகிறோம்.
\end{proof}

% SEGMENT OLP-0035-S08
\begin{prob}
$\cardeq{A}{C}$, $\cardeq{B}{D}$ எனவும்,
$A \cap B =
C \cap D = \emptyset$ எனவும் இருந்தால்,
$\cardeq{A \cup B}{C \cup D}$ எனக் காட்டுக.
\end{prob}

% SEGMENT OLP-0035-S09
\begin{prob}
$A$ முடிவுறாததாகவும் !!{enumerable} ஆகவும் இருந்தால்,
$\cardeq{A}{\Nat}$ எனக் காட்டுக.
\end{prob}

\end{document}
